%!TEX root = ../../exa-ma-d7.3.tex
\subsection{Software: Composyx}
\label{sec:Composyx:software:training}



\begin{table}[H]
    \centering
    { \setlength{\parindent}{0pt}
    \def\arraystretch{1.25}
    \arrayrulecolor{numpexgray}
    {\fontsize{9}{11}\selectfont
    \begin{tabular}{!{\color{numpexgray}\vrule}p{.4\textwidth}!{\color{numpexgray}\vrule}p{.6\textwidth}!{\color{numpexgray}\vrule}}
        \rowcolor{numpexgray}{\rule{0pt}{2.5ex}\color{white}\bf Field} & {\rule{0pt}{2.5ex}\color{white}\bf Details} \\
        \rowcolor{white}\textbf{Consortium} & \begin{tabular}{l}
None\\
\end{tabular} \\
        \rowcolor{numpexlightergray}\textbf{Exa-MA Partners} & \begin{tabular}{l}
Inria BXSO\\
\end{tabular} \\
        \rowcolor{white}\textbf{Contact Emails} & \begin{tabular}{l}
gilles.marait@inria.fr\\
\end{tabular} \\
        \rowcolor{numpexlightergray}\textbf{Supported Architectures} & \begin{tabular}{l}
CPU Only\\
\end{tabular} \\
        \rowcolor{white}\textbf{Repository} & \href{https://gitlab.inria.fr/composyx/composyx}{https://gitlab.inria.fr/composyx/composyx} \\
        \rowcolor{numpexlightergray}\textbf{License} & \begin{tabular}{l}
OSS: Cecill-*\\
\end{tabular} \\
        \rowcolor{white}\textbf{Training} & \begin{tabular}{l}
          Documentation \\
          Tutorials \\
\end{tabular} \\
        \bottomrule
    \end{tabular}
    
%%\begin{table}[!ht]
%%    \centering
    {
    \setlength{\parindent}{0pt}
    \def\arraystretch{1.25}
    \arrayrulecolor{numpexgray}
    {\fontsize{9}{11}\selectfont
    \begin{tabular}{!{\color{numpexgray}\vrule}p{0.20\textwidth}
                    !{\color{numpexgray}\vrule}p{0.40\textwidth}
                    !{\color{numpexgray}\vrule}p{0.35\textwidth}
                    !{\color{numpexgray}\vrule}}
    \rowcolor{numpexgray}{\rule{0pt}{2.5ex}\color{white}\bf Training Category} 
         & {\rule{0pt}{2.5ex}\color{white}\bf Description} 
      & {\rule{0pt}{2.5ex}\color{white}\bf URL/Links} \\

      \rowcolor{white}Documentation & Full documentation (html pages) & \url{https://composyx.gitlabpages.inria.fr/composyx/main/} \\

      \rowcolor{numpexlightergray} Documentation & Full documentation (pdf file) & \url{https://composyx.gitlabpages.inria.fr/composyx/main/documentation_full.pdf} \\

      \rowcolor{white} Tutorial & Hand-on tutorial & \url{https://composyx.gitlabpages.inria.fr/composyx/main/tutorial/tuto.html} \\

      \rowcolor{numpexlightergray} Tutorials & Usage and parameters of GMRES in Composyx & \url{https://composyx.gitlabpages.inria.fr/composyx/main/tutorial/gmres.html} \\

      \rowcolor{white} Tutorial & Definition of customized operators in Composyx & \url{https://composyx.gitlabpages.inria.fr/composyx/main/tutorial/operators.html} \\

      \rowcolor{numpexlightergray} Tutorials & Using C and fortran interfaces & \url{https://composyx.gitlabpages.inria.fr/composyx/main/tutorial/f_c_drivers.html} \\

\bottomrule
    \end{tabular}
    }}
%    \caption{Composyx Training}
%\end{table}

    }}
    \caption{Composyx Information}
\end{table}


\paragraph{Additional Details:}
\begin{description}
\item[\textbf{Documentation}] % start of item
\emph{Full documentation for users and developers.}

\begin{itemize}
\item \textbf{Topics Covered:} The whole code of Composyx is annoted
  and commented in litterate programming.
  \item \textbf{Format:} The code is parsed to generate the documentation in \textit{html} pages and a \textit{pdf} file.
  \item \textbf{Timeline:} In literate programming, the documentation
    progresses hand-to-hand with the code, although some factorization
    efforts are sometimes needed to keep the documentation clear and
    complete.
\end{itemize}

\item[\textbf{Tutorials}] % start of item
\emph{Hands-on or step-by-step guides for newcomers and practitioners, and tutorials for more specific features.}

\begin{itemize}
\item \textbf{Modules Included:} Software deployment, basic concepts, up to a parallel MPI example
\item \textbf{Target Audience:} Any user of Composyx
\item \textbf{Existing Materials:} \begin{itemize}
  \item \url{https://composyx.gitlabpages.inria.fr/composyx/main/tutorial/tuto.html}
  \item \url{https://composyx.gitlabpages.inria.fr/composyx/main/tutorial/gmres.html}
  \item \url{https://composyx.gitlabpages.inria.fr/composyx/main/tutorial/operators.html}
  \item \url{https://composyx.gitlabpages.inria.fr/composyx/main/tutorial/f_c_drivers.html}
  \end{itemize}
\end{itemize}

\end{description}

\subsubsection{Next Steps for Composyx}
\label{sec:Composyx:next-steps}

The Composyx software is growing large and it has been noted that
users and developers have difficulties finding features and interfaces
which are effectively available in the software. Some efforts to make
reference cards and links between parts of the documentation are
required.

\begin{itemize}
\item \textbf{New Training Content:} Upgrade of the documentation, to
  make it easier to find available features. New tutorial on specific
  topics.
\item \textbf{Incorporating Exascale Techniques and User Feedback:}
  New benchmarks, with documentation on how to produce the results
  will be added to the documentation.
    \item \textbf{Delivery Formats:} At the moment, only online
      documentation and \textit{pdf} files are considered.
    \item \textbf{Enhancing Support Tools:} Addition of new
      topic-related tutorials and benchmark documentation.
    \item \textbf{Integrating Training with Academic and Industrial
        Outreach:} Composyx is developed in the the concace project is
      a joint Inria-Industry team involving members from two private
      partners, namely, Airbus Central R\&T and Cerfacs. Hopefully the
      integration of Composyx in other academic and industrial
      software in the future, from Inria, Airbus, Cerfacs or other
      collaborations shall create opportunities to improve the
      training material of Composyx.
\end{itemize}
